% The Zephyr build, in the order it happens: what each step reads, what it writes, and where the
% two phases divide. Drawn to match §1.2's "What happens when", so the figure and the text name
% the same files.
%
% The inks say what kind of thing something is, as they do in the code blocks: green for what the
% build is configured with (devicetree and Kconfig), blue for the build system itself, red for C
% and the program it becomes.
\begin{tikzpicture}[
    x=1cm, y=1cm, font=\footnotesize,
    stage/.style={draw=#1, fill=fill#1, line width=0.7pt, rounded corners=1.5pt,
      minimum width=3.15cm, minimum height=0.78cm, align=center, font=\small\bfseries,
      text=#1!60!black},
    artifact/.style={align=left, anchor=west, font=\scriptsize\ttfamily, text=ink},
    input/.style={align=right, anchor=east, font=\scriptsize\ttfamily, text=muted},
    flow/.style={-{Latex[length=2mm]}, line width=0.9pt, ink},
    feed/.style={-{Latex[length=1.6mm]}, line width=0.5pt, muted},
    writes/.style={-{Latex[length=1.6mm]}, line width=0.5pt, rule},
  ]
  % One row per step, top to bottom.
  \def\x{0}
  \def\ya{0} \def\yb{-1.55} \def\yc{-3.1} \def\yd{-4.65} \def\ye{-6.5}

  \node[stage=muted] (board) at (\x, \ya) {1\enspace Board};
  \node[stage=accent3] (dt) at (\x, \yb) {2\enspace Devicetree};
  \node[stage=accent3] (kconfig) at (\x, \yc) {3\enspace Kconfig};
  \node[stage=accent2] (cmake) at (\x, \yd) {4\enspace CMake};
  \node[stage=accent] (ninja) at (\x, \ye) {5\enspace Ninja};

  \draw[flow] (board) -- (dt);
  \draw[flow] (dt) -- (kconfig);
  \draw[flow] (kconfig) -- (cmake);
  \draw[flow] (cmake) -- (ninja);

  % What each step reads, on the left.
  \node[input] (iboard) at (\x-2.2, \ya) {-b native\_sim/native/64};
  \node[input] (idt) at (\x-2.2, \yb) {native\_sim.dts, app.overlay};
  \node[input] (ikconfig) at (\x-2.2, \yc) {Kconfig, prj.conf,\\-DCONFIG\_...};
  \node[input] (icmake) at (\x-2.2, \yd) {CMakeLists.txt};
  \node[input] (ininja) at (\x-2.2, \ye) {src/*.c, Zephyr's sources};
  \foreach \n in {iboard, idt, ikconfig, icmake, ininja} {
    \draw[feed] (\n.east) -- ++(0.35, 0);
  }

  % What each step writes, on the right.
  \node[artifact] (aboard) at (\x+2.2, \ya) {the SoC, the arch,\\the board's defaults};
  \node[artifact] (adt) at (\x+2.2, \yb) {build/zephyr/zephyr.dts\\devicetree\_generated.h};
  \node[artifact] (akconfig) at (\x+2.2, \yc) {build/zephyr/.config\\autoconf.h};
  \node[artifact] (acmake) at (\x+2.2, \yd) {build/build.ninja};
  \node[artifact, text=accent!70!black] (aninja) at (\x+2.2, \ye)
    {build/zephyr/zephyr.exe};
  \foreach \n in {aboard, adt, akconfig, acmake, aninja} {
    \draw[writes] (\n.west) -- ++(-0.35, 0);
  }

  % The two phases, named beside a rule, which costs less width than a brace and its label.
  \draw[rule, line width=0.6pt] (\x-5.75, \ya+0.45) -- (\x-5.75, \yd-0.45);
  \node[rotate=90, anchor=south, font=\scriptsize\scshape\lsstyle, text=ink]
    at (\x-5.8, {(\ya+\yd)/2}) {configuration};
  \draw[rule, line width=0.6pt] (\x-5.75, \ye+0.45) -- (\x-5.75, \ye-0.45);
  \node[rotate=90, anchor=south, font=\scriptsize\scshape\lsstyle, text=ink]
    at (\x-5.8, \ye) {compilation};
\end{tikzpicture}
